% short/05_methods.tex: how the models were tested (LOSO, and the training-only split, which gives predictions for the
% training participants from models not trained on them, with the reason), the measures (accuracy, AUC, log loss =
% binary cross-entropy, 95 % person-bootstrap CI), the six models in three named families plus a baseline (MiniRocket
% and the CNN explained in plain words), mirror augmentation (after the models, which it refers to), calibration
% (Platt), combining attempts (intuition and a worked example, then equation eq:aggregate), how the processing choices
% were tested (two fixed test models, MiniRocket for several options; clearly worse in accuracy and in log loss), and
% the choice of the final model with its honest (repeated-choice) estimate. The model's score is written s (not z,
% which names the phone's axis out of the screen). Numbers only through \R keys.
\section{Models and how we tested them}\label{sec:methods}

\paragraph{Testing on new people.} To measure how a model does on a person it has never seen, we used
leave-one-subject-out (LOSO) cross-validation: a model is trained on \R{final/n_train_persons} participants and tested
on the remaining one, and this is repeated \R{describe/n_persons} times, so that every participant is the test
participant once. Everything a model learns from data, including its calibration and the factor used to combine
attempts (both described below), is learned from the \R{final/n_train_persons} training participants only. Some of
these steps need predictions for the training participants made by a model that was not trained on them (a model
predicts the people it was trained on better than new people, so calibrating or combining on such predictions would
make it over-confident). For these steps we split the training participants into a few groups, trained on all groups
but one and predicted the group left out, in turn; we call this the \emph{training-only split}. The predictions of the
\R{describe/n_persons} LOSO rounds are then pooled and scored together.

\paragraph{Measures.} The \emph{accuracy} is the share of trials in which the hand given more than 50\,\% is the hand
actually used. The AUC of Section~\ref{sec:physics} measures how well the predicted \pright{} separates the hands. The
\emph{log loss}, or binary cross-entropy, is the average over the trials of $-\ln$ of the probability given to the hand
actually used: it is low when the model is right and confident and punishes confident mistakes heavily; a model that
always answers 50\,\% scores \R{final/log_loss_chance}, and lower is better. The main numbers come with a
95\,\% confidence interval (CI), written [low, high]: unless stated otherwise, we resampled the \R{describe/n_persons} participants with
replacement \R{p1/n_boot} times, recomputed the number each time and kept the middle 95\,\% of these values (a person
bootstrap, which respects that the trials of one person are alike).

\paragraph{The models.} Because one number is not enough, we built models that read more of the movement, in three
families: a physics model that uses a few numbers chosen from the mechanism of Section~\ref{sec:physics}; standard
machine learning on many summary numbers of each trial (gradient boosting and a stack of three models); and models
that read the whole time series (MiniRocket and a 1D CNN). A simple yaw baseline completes the comparison. Every model
reads the movement to the nose of Section~\ref{sec:processing} and outputs a score $s$, the log-odds of the right hand,
$s = \ln\bigl(\pright/(1-\pright)\bigr)$: $s = 0$ means 50\,\%, positive values favour the right hand, and
$\pright = \sigma(s) = 1/(1 + e^{-s})$.
\begin{itemize}[nosep,leftmargin=1.2em]
  \item \emph{Yaw baseline.} A logistic regression on one channel, the yaw, read at \R{final/ref_n_instants} evenly
    spaced moments of the movement. It was fixed before any other model was built, as the baseline that a more
    complex model has to beat.
  \item \emph{Physics model.} An antisymmetric logistic regression on \R{final/phys6_n_features} numbers chosen from
    physics: the yaw, the roll (the turn about the towards-the-face direction so far) and the sideways velocity (the
    sideways acceleration added up), each read halfway through the movement and at its end. All six change sign
    under the mirror and the regression has no intercept, so the mirror flips the sign of $s$, and since
    $\sigma(-s) = 1 - \sigma(s)$ the mirror image of a movement gets exactly $1 - \pright$: the two hands are
    treated symmetrically by construction.
  \item \emph{Gradient boosting} (LightGBM): \R{final/gbm_n_trees} small decision trees, each correcting the errors of
    the previous ones. A decision tree sorts a trial by a few yes/no questions about its numbers (for example, is the
    average yaw in some time segment above a threshold?), and each final group, a leaf (at most \R{final/gbm_n_leaves}
    per tree), adds its own amount to the score. The trees read \R{final/gbm_n_features} summary
    numbers: for each of \R{final/gbm_n_channels} channels (the \R{final/cnn_n_gravity} gravity-frame channels and,
    along each of the phone's three axes, the acceleration, the rotation speed and the direction of up, which shows
    how the phone is held), its average in \R{p1/n_bins} equal time segments, its mean, spread, minimum and maximum,
    and when these two occur.
  \item \emph{Stack of \R{final/stack_n_models} models.} A logistic regression, a support vector machine (a classifier
    that draws a curved boundary between the hands) and gradient boosting, each on \R{m2/lr_full_n_features} summary
    numbers per trial (averages in time segments, statistics, frequency content, correlations between channels and
    running sums, among others), combined by a small logistic regression that learns how much to trust each, from
    their predictions in the training-only split.
  \item \emph{MiniRocket}~\cite{dempster2021minirocket}, a fast general-purpose method for time series. It slides a
    fixed set of short patterns along the \R{final/cnn_n_gravity} gravity-frame channels: at each moment a pattern is
    multiplied point by point with the stretch of signal under it and the products are added up, which gives a large
    response where the signal looks like the pattern. Each pattern is also stretched to several time scales, and for
    each the share of the movement where it responds becomes one summary number (\R{final/mr_n_features_max} in all);
    a logistic regression is fitted on these numbers.
  \item \emph{1D CNN}, the final model: a small one-dimensional convolutional neural network that reads the
    \R{final/cnn_n_channels} gravity- and body-frame channels as time series of \R{describe/n_norm} points. Like
    MiniRocket it slides short patterns along the signals, but it learns them from the training data, and it stacks
    them in \R{final/cnn_n_layers} layers of \R{final/cnn_width} patterns, \R{final/cnn_kernels} points long: the
    first layer compares its patterns with the channels, and each later layer looks for patterns in the responses of
    the layer before, so that it finds longer patterns built from shorter ones. The responses are averaged within
    \R{final/cnn_n_bins} time segments, so that the network keeps when in the movement a pattern occurred, and a
    weighted sum of these averages gives $s$. The numbers in the patterns and in the sum, the weights
    (\R{final/cnn_n_params} per network), are learned in many small steps that lower the log loss on the training
    trials. We split the training participants into \R{final/cnn_n_members} groups and
    trained \R{final/cnn_n_members} such networks, each on all groups but one, stopping when its log loss on the group
    left out no longer improved (which makes its predictions for that group slightly flattering); for a new
    participant their outputs are averaged. Training also used small random stretches in time, changes of scale and
    noise.
\end{itemize}

\paragraph{Mirror images as extra training data.} Every model but the physics model, which is symmetric by
construction, was also trained on the mirror image of every training movement, labelled with the other hand (mirror
augmentation): the same channels with every sideways motion, every turn about the vertical and every roll reversed
(Section~\ref{sec:physics}). This doubles the training data, balances the two hands and teaches the model that the two
hands move as mirror images.

\paragraph{Calibration.} A model's output need not be a calibrated probability: many models are over-confident, with
probabilities too close to 0 and 1~\cite{guo2017calibration}. We check calibration with a reliability curve: the trials are
sorted by their predicted \pright{} and cut into \R{final/calfig_cnn_groups} groups of \R{final/calfig_cnn_group_min}
or \R{final/calfig_cnn_group_max} trials, and for each group the average predicted \pright{} is plotted against the
share of its trials that really were right-hand trials; a calibrated model lies on the diagonal. A calibration function
can correct a model that is not calibrated. We used Platt scaling~\cite{platt1999probabilistic}: the calibrated
probability is $\sigma(a s + b)$, where $s$ is the model's score; $a < 1$ pulls extreme scores towards 50\,\%,
which corrects over-confidence, $a > 1$ pushes them away from 50\,\%, which corrects under-confidence, and $b$ shifts
all scores towards one hand. In each LOSO round, $a$ and $b$ are fitted on the training participants' predictions from
the training-only split, so the left-out participant never informs its own calibration.

\paragraph{Combining several attempts.} A person can repeat the movement with the same hand, and each such attempt
(one more kept trial of the same person and hand) adds evidence about the hand. On the log-odds scale the evidence of
independent observations simply adds up: two independent answers of 80\,\% each are odds of 4 to 1 each, together
$4 \times 4 = 16$ to 1, that is 94\,\%, and multiplying odds is adding log-odds. The attempts of one person are not
independent, however (same person, same habits), so adding them as if they were would make the combined answer
over-confident. We therefore add the scores of the $k$ attempts, divide the sum by a factor \tauk{}, and turn the
result back into a probability:
\begin{equation}
  \pright_k = \sigma\!\left(\frac{1}{\tauk}\sum_{i=1}^{k} s_i\right),
  \label{eq:aggregate}
\end{equation}
where $s_i$ is the score of attempt $i$. The factor \tauk{} lies between 1 (attempts fully independent: a plain sum)
and $k$ (attempts identical: a plain average of the scores). In each LOSO round it is the value that gives the lowest
log loss on random groups of $k$ attempts of the training participants, scored with their predictions from the
training-only split. For three attempts its median over the rounds is \R{final/att_tau3}: dividing the sum of three
scores by \R{final/att_tau3} counts them as \R{final/att_mult3} independent attempts instead of 3, so three attempts
are worth about two independent ones.

\paragraph{How we chose the processing.} For the processing steps of Section~\ref{sec:processing}, except the iPhone
sign (a correction, not a choice) and the body frame (which only the CNN reads), we also tried alternatives, re-running
the whole LOSO evaluation each time with two fixed test models (a logistic regression on the yaw rate and the yaw,
similar to the yaw baseline, and gradient boosting) and, for several alternatives, also with MiniRocket: no smoothing
or another filter, a coarser clock, the phone's own estimate of gravity, a window of fixed length instead of the
time-normalised movement, channels rescaled to a common size, other sets of channels and no mirror images, among
others (\R{p1/n_choices} choices in all). No alternative was clearly better. In accuracy, two were clearly worse:
replacing the yaw rate and the yaw by the signals along the phone's own axes changed the accuracy of the logistic
regression by \Rcib{p1/d_lr_frame_device} percentage points, and keeping only the gravity-frame channels changed that
of gradient boosting by \Rcib{p1/d_gbm_frame_gravity} points (the phone's own axes still carry some information for
this model; the final model nevertheless reads only channels that do not change with the grip,
Section~\ref{sec:processing}). By log loss, a few more were clearly worse for gradient boosting, among them training
without the mirror images (a change of \Rcib{p1/dll_gbm_no_mirror}). All other differences were of a size that chance
alone can produce. We therefore kept the processing described above.

\paragraph{Choosing the final model, and an honest estimate.} We compared \R{final/n_candidates} models in all
(several versions within each family) and kept the 1D CNN, which had the lowest log loss and the highest AUC. Picking
the best of many models on the same participants makes the winner's own score slightly flattering, because part of its
lead can be luck. To get an honest number for a new participant we repeated the choice inside every LOSO round: with
one participant left out, the \R{final/n_candidates} models were compared by their log loss on the other
\R{final/n_train_persons} participants only (with the training-only split, in \R{final/nested_inner_folds} groups),
and the one with the lowest log loss there was tested on the left-out participant. In a round the choice can fall on
another model, so this estimate measures our whole way of choosing a model, applied to a new person; on all
\R{describe/n_persons} participants the same choice gives the CNN. This honest estimate is our main result. It repeats
the final choice among the models, but not the earlier design decisions made while looking at the same results (for
example, which channels the CNN reads), so it may still be slightly optimistic.
